% auto-generated by consolidated_report.ipynb - do not edit by hand
\begin{table}[htbp]
\centering
\caption{Open items blocking the transfer to LaTeX}
\label{tab:note_1_open_items}
\small
\begin{tabular}{>{\raggedright\arraybackslash}p{\dimexpr 0.0489\textwidth-2\tabcolsep\relax}>{\raggedright\arraybackslash}p{\dimexpr 0.2344\textwidth-2\tabcolsep\relax}>{\raggedright\arraybackslash}p{\dimexpr 0.3167\textwidth-2\tabcolsep\relax}>{\raggedright\arraybackslash}p{\dimexpr 0.4000\textwidth-2\tabcolsep\relax}}
\toprule
\# & Open item & Evidence & Effect on the paper \\
\midrule
N1 & Japan's study window is two formation years (2023, 2024) & japan\_\allowbreak{}diagnostics.csv; japan\_\allowbreak{}fullperiod\_\allowbreak{}feasibility.csv reports “no point-in-time statements” for 2005--2022 & No 13-year window is common to all three markets. §8's headline cross-country table cannot be described as covering 2012--2024. \\
N2 & Japan's 2023 basket holds 9 names, not K & japan\_\allowbreak{}diagnostics.csv (k = 9 in 2023, 30 in 2024) & Half of Japan's headline result rests on a 9-name portfolio. The Effective N of 19.5 is the average of the two years, not a basket size. \\
N3 & The eligible universe is capped at the top 150 names by dollar volume & pipeline.py:105,141; the cap binds in every market-year except Japan 2023 (118) and US 2010 (121) & After the 40\% B/\allowbreak{}M cut the null draws from 60 names, so a random basket of K = 30 shares half its names with the F-Score basket. The Monte Carlo test has little power by construction. \\
N4 & The high-B/\allowbreak{}M set is 40\% of the cap, not 40\% of the market & high\_\allowbreak{}bm\_\allowbreak{}subset takes a relative quantile (loaders.py:78); value\_\allowbreak{}set = 60 in every year where universe = 150, without exception & The count cannot respond to the market: a value drawdown does not widen it. Uncapped, 40\% would be 121--342 names in Vietnam. The B/\allowbreak{}M threshold is also a moving target, differing by year and by market. \\
N5 & Monte Carlo N differs by strategy (1000 EW /\allowbreak{} 300 GMV) & *\_\allowbreak{}mc\_\allowbreak{}placement.csv, n\_\allowbreak{}draws column & §6 requires one N for all three headline country tests. GMV p-values are quoted on a coarser grid than EW p-values. \\
N6 & The US and Japan fundamentals caches are absent & data/\allowbreak{} holds only vietnam\_\allowbreak{}*; loaders.py:47 reads data/\allowbreak{}{market}\_\allowbreak{}fundamentals.csv for both markets & The US and Japan headline results cannot be regenerated as they stand. scripts/\allowbreak{}fetch\_\allowbreak{}us\_\allowbreak{}japan.py must rebuild the caches first, and refreshed Yahoo data may not reproduce the committed numbers. \\
N7 & Maximum weight and non-zero name counts are not exported & *\_\allowbreak{}summary.csv carries effective\_\allowbreak{}n only & P24 asks for both. Those two columns stay as --- in Table 8.1 until the pipeline writes them. \\
N8 & Drawdown and annual-weight figures are not produced & results/\allowbreak{}figures/\allowbreak{} holds NAV, MC placement and F-Score distribution only & P17, P20 and P23 can be filled only in part; see the figure audit below. \\
N9 & The $-100\%$ delisting robustness case was not run & no results file for it & P28 must either be run or state plainly that it was not. \\
N10 & Sector-capped GMV exists in all three markets but is placed nowhere & fscore\_\allowbreak{}GMVsec rows in every *\_\allowbreak{}summary.csv & P10 and P29 both wait on this decision: headline or robustness. \\
\bottomrule
\end{tabular}
\begin{minipage}{\linewidth}\vspace{4pt}\footnotesize
Compiled from the results files and the implementation, not from the draft. N1--N5 constrain what the numbers may be said to show; N6--N10 are gaps in what has been produced.
\end{minipage}
\end{table}
